<!DOCTYPE html>
<!-- Generated by pkgdown: do not edit by hand --><html lang="en"><head><meta http-equiv="Content-Type" content="text/html; charset=UTF-8"><meta charset="utf-8"><meta http-equiv="X-UA-Compatible" content="IE=edge"><meta name="viewport" content="width=device-width, initial-scale=1, shrink-to-fit=no"><title>BIOL 322 Lab 5. Gene Flow • biol322</title><script src="../deps/jquery-3.6.0/jquery-3.6.0.min.js"></script><meta name="viewport" content="width=device-width, initial-scale=1, shrink-to-fit=no"><link href="../deps/bootstrap-5.3.1/bootstrap.min.css" rel="stylesheet"><script src="../deps/bootstrap-5.3.1/bootstrap.bundle.min.js"></script><link href="../deps/font-awesome-6.5.2/css/all.min.css" rel="stylesheet"><link href="../deps/font-awesome-6.5.2/css/v4-shims.min.css" rel="stylesheet"><script src="../deps/headroom-0.11.0/headroom.min.js"></script><script src="../deps/headroom-0.11.0/jQuery.headroom.min.js"></script><script src="../deps/bootstrap-toc-1.0.1/bootstrap-toc.min.js"></script><script src="../deps/clipboard.js-2.0.11/clipboard.min.js"></script><script src="../deps/search-1.0.0/autocomplete.jquery.min.js"></script><script src="../deps/search-1.0.0/fuse.min.js"></script><script src="../deps/search-1.0.0/mark.min.js"></script><!-- pkgdown --><script src="../pkgdown.js"></script><link href="../extra.css" rel="stylesheet"><meta property="og:title" content="BIOL 322 Lab 5. Gene Flow"></head><body>
    <a href="#main" class="visually-hidden-focusable">Skip to contents</a>


    <nav class="navbar navbar-expand-lg fixed-top bg-primary" data-bs-theme="dark" aria-label="Site navigation"><div class="container">

    <a class="navbar-brand me-2" href="../index.html">biol322</a>

    <small class="nav-text text-muted me-auto" data-bs-toggle="tooltip" data-bs-placement="bottom" title="">1.0</small>


    <button class="navbar-toggler" type="button" data-bs-toggle="collapse" data-bs-target="#navbar" aria-controls="navbar" aria-expanded="false" aria-label="Toggle navigation">
      <span class="navbar-toggler-icon"></span>
    </button>

    <div id="navbar" class="collapse navbar-collapse ms-3">
      <ul class="navbar-nav me-auto"><li class="nav-item"><a class="nav-link" href="../index.html"><span class="fa fa-home fa-lg"></span></a></li>
<li class="nav-item dropdown">
  <button class="nav-link dropdown-toggle" type="button" id="dropdown-course-materials" data-bs-toggle="dropdown" aria-expanded="false" aria-haspopup="true">Course Materials</button>
  <ul class="dropdown-menu" aria-labelledby="dropdown-course-materials"><li><hr class="dropdown-divider"></li>
    <li><a class="dropdown-item" href="../articles/sheets/sci_articles/biol322-syllabus-Roles-spring2026.pdf">Syllabus PDF</a></li>
    <li><a class="external-link dropdown-item" href="https://docs.google.com/document/d/1sV-864Y8L7gAb3LwBSqtKKL0P0aBKue20PB63_IPu-Y/edit?usp=sharing">Grade Contract Template</a></li>
    <li><a class="external-link dropdown-item" href="https://drive.google.com/drive/folders/1arcg7_nXRov4kTE3jgnS3q8DQ9DwZcFO?usp=sharing">Course GDrive folder</a></li>
    <li><a class="external-link dropdown-item" href="https://drive.google.com/drive/folders/1xDjUUTZY1j88zHregL34JMnaW6BZAUzq?usp=drive_link">In-class materials</a></li>
    <li><a class="dropdown-item" href="../articles/lab-exercises-main.html">Lab Exercises</a></li>
    <li><a class="external-link dropdown-item" href="https://drive.google.com/drive/folders/1ZuXCOFdWGMYW6JMkuYXSKr1I2bELQYDm?usp=drive_link">Exams</a></li>
    <li><a class="external-link dropdown-item" href="https://docs.google.com/document/d/1VkC-BHN8_EGvSv86qwt3uVwcK--VN61rOL4qlOGoAos/edit?usp=drive_link">Midterm Reflection</a></li>
    <li><a class="external-link dropdown-item" href="https://drive.google.com/drive/folders/18AjoufAqq3saqJmHq8bZrmXG0pGoEM00?usp=drive_link">Student Folders</a></li>
    <li><a class="external-link dropdown-item" href="https://docs.google.com/spreadsheets/d/1cfcKbMkjukYi1RewOfqhoLowUIsxJcsT34LbCcgj39I/edit?usp=drive_link">Check Progress Spreadsheet</a></li>
    <li><hr class="dropdown-divider"></li>
  </ul></li>
<li class="nav-item"><a class="external-link nav-link" href="https://docs.google.com/document/d/1gNcEPhAxowrrbwCOjgP2J_djOn_QOkaKBlhnRhSO5iQ/edit?usp=drive_link">Course Agenda</a></li>
<li class="nav-item"><a class="external-link nav-link" href="https://docs.google.com/forms/d/e/1FAIpQLSdA2mkUojv-37gaK0ON-B8B2a7y5cxnMUJRhaqX747QtbMxhA/viewform?usp=header">Assignment Turn In</a></li>
      </ul><ul class="navbar-nav"><li class="nav-item"><a class="nav-link" href="../articles/sheets/sci_articles/biol322-syllabus-Roles-spring2026.pdf">BIOL 322 Spring 2026 Syllabus PDF</a></li>
      </ul></div>


  </div>
</nav><div class="container template-article">


<div class="row">
  <main id="main" class="col-md-9"><div class="page-header">

      <h1>BIOL 322 Lab 5. Gene Flow</h1>
                        <h4 data-toc-skip class="author">Angela
Roles</h4>
            
            <h4 data-toc-skip class="date">Spring 2026</h4>
      

      <div class="d-none name"><code>biol322-lab5.Rmd</code></div>
    </div>

    
    
When you have completed the exercise below, submit your model and your
lab write-up via the \href{https://forms.gle/E3QLACX1aA2Q6ist5}{Google
Drive Form}. Responses due by the following week's lab period.

During lab, we will work in groups. Here are the groups assigned for
today's exercise:

\begin{itemize}
\tightlist
\item
  Savannah and Natalie
\item
  Charlie and Paige
\item
  Ella and Arc
\item
  Amanda, Clara, and Jordyn
\item
  Emmy and Amy
\item
  Miles and Bethany
\end{itemize}

\begin{center}\rule{0.5\linewidth}{0.5pt}\end{center}

\subsection{Overview}\label{overview}

Today, we're ready to add gene flow to our model, alongside selection
and drift. After completing today's model, we'll have a creation that
allows us to investigate the effects of the three major forces of
evolutionary change that we can observe on a human timescale! And, we'll
have set it up so we can allow for any combination of the three
occurring at once (e.g., gene flow alone versus selection and drift). We
have only just covered this chapter in class. So let me guide you
through how to think about adding gene flow to our model.

In the real world, we might often be thinking about gene flow occurring
as individuals migrate from one sub-population into another
sub-population. However, we could also focus on a single sub-population
of interest and only consider how migration is influencing allele and
genotype frequencies in that single sub-population. For simplicity,
that's what we'll do with our model today. There are also ways to create
multiple instances of our basic model and enable the models to interact
with each other but for now, we'll stick with the single sub-population.
After all, this is how we think about the island-continent model of gene
flow that we're working on today. In our case, the island is our
population and the continent is where migrating individuals (or gametes)
come from. In this particular model, gene flow is only coming from the
continent to the island, so if only gene flow is occurring, we always
predict that the island allele frequency will eventually be the same as
that on the continent. Thus, the continent allele frequency is our gene
flow equilibrium.

In order to add gene flow to our existing model, we'll presume that new
alleles are added to the gene pool (read: the pool of gametes), either
the production of gametes by newly arrived adults or via the movement of
gametes themselves (such as pollen). To do this, we will adjust the
allele frequency in our gametes, before applying drift to those allele
frequencies. And then selection will still occur on the viability of
zygotes (survival to maturity).

Assuming we're intending to model the island-continent version of gene
flow, that means we have just 2 variables that need to be added to our
current model:

\begin{itemize}
\tightlist
\item
  pm, the frequency of allele A1 in the continent population (where our
  migrants are coming from; range 0-1).
\item
  m, the rate of migration as a proportion; that is, the fraction of
  alleles in our sub-population coming from individuals not born in this
  sub-population (range 0-1).
\end{itemize}

Given a starting allele frequency of pt in our sub-population, here's
our equation for allele frequency in the next generation of our
sub-population:

pt+1 = (1−m) * pt + m * pm

That means we'll adjust the existing variable p, updating it's equation
to match the one above. At present, our variable p is defined as, p =
freqA1A1 + 0.5 ∗ freqA1A2. When we update this equation, we use the
existing value (p) for pt in the equation including gene flow. Note that
when m = 0, that will be the ``no gene flow'' scenario so that's how to
turn gene flow `on' and `off'.

Here's a reference image of the completed model. Compared to last week,
the only visible difference is the addition of those two components (m
and pm) plus one extra, called selecteq which will be described below.

\includegraphics[alt={Image of the completed model with selection, drift, and gene flow}]{../man/figures/geneflowmodel}

\begin{center}\rule{0.5\linewidth}{0.5pt}\end{center}

\subsection{Add Gene Flow to the Model}\label{addflow}

As described above, we need to add a rate of migration, m, and the
frequency of allele A1 in the source population, pm, to our model. Both
of these components will be variables we adjust the value of so we'll go
ahead and specify a min, max, and increment in addition to a starting
value.

\begin{enumerate}
\def\labelenumi{\arabic{enumi}.}
\item
  Create these two variables then click Equation and set them up.

  Variable

  Description

  Equation

  Min

  Max

  Incr

  pm

  frequency of A1 allele in migrants

  0

  0

  1

  0.001

  m

  rate of migration, proportion of sub-pop alleles from individuals not
  born in the sub-pop

  0

  0

  1

  0.001

  Keep in mind that a value of m = 0.1 means that each generation 10\%
  of the alleles in the population are from migrants. The actual number
  of migrant individuals arriving in the population will be the
  migration rate times the total population size (Nm).
\item
  Now, we need to use these variables in the rest of the model.
  Specifically, as described, we'll adjust the value of p using the
  equation from above. We'll multiply the existing value of p by (1 − m)
  and then add mpm to the whole thing. This means that in the life cycle
  of our organisms, we are allowing gene flow to happen prior to gamete
  production, then drift occurs to the gene pool, adjusting allele
  frequencies (yielding p adj), then zygotes are created (following HW)
  and their survival to maturity (viability) is adjusted by selection in
  the flows. Here's the equation for p in our gene flow model: (freqA1A1
  + 0.5 * freqA1A2) * (1-m) + pm*m Since pm is our equilibrium allele
  frequency for gene flow, we can compare our actual allele frequency to
  see whether or not we reach equilibrium. Let's add one more Variable
  that will let us display the selection-only equilibrium too.
\item
  Create a variable called selecteq. This will take the settings we
  chose for fitness and determine what value of p we expect at
  equilibrium. To do that, we'll set the equation as an if-then
  statement, like we used in the drift model in Lab 3: IF THEN
  ELSE(wA1A1 = wA1A2 :AND: wA1A1 = wA2A2, p0, IF THEN ELSE((wA1A2
  \textgreater{} wA1A1 :AND: wA1A2 \textgreater{} wA2A2) :OR: (wA1A2
  \textless{} wA2A2 :AND: wA1A2 \textless{} wA1A1), (wA1A2 - wA2A2) / (2
  * wA1A2 - wA1A1 - wA2A2), IF THEN ELSE(wA1A1 = 1, 1, 0))) This is a
  super complicated equation so be careful to enter it correctly. If you
  copy and paste it should be ok but please try to be sure you
  understand what the equation means\ldots{}

  \begin{itemize}
  \tightlist
  \item
    In short, it says that IF all of the genotypes have the same
    fitness, THEN the equilibrium value will be whatever p is to begin
    with, p0. ELSE, if genotype fitnesses are not equal, THEN, IF A1A2
    has the highest OR lowest fitness, THEN use the equation for
    heterozygote advantage to figure out the equilibrium. ELSE, IF one
    of the homozygotes has the highest fitness, THEN IF A1A1 has the
    highest fitness set equilibrium to p=1.0, ELSE set equilibrium to
    p=0 (because A2A2 has the highest fitness and will reach fixation).
  \end{itemize}
\item
  Now, when we are modeling gene flow, we'll want to graph p, pm, and
  selecteq so we can see how the actual outcome compares to the
  equilibrium we'd predict if only gene flow were happening (pm) or only
  selection were happening (selecteq). We're now ready to do a test run
  before using the model to understand gene flow and how it interacts
  with drift and selection.
\end{enumerate}

\subsection{Test the model}\label{test-the-model}

\begin{enumerate}
\def\labelenumi{\arabic{enumi}.}
\setcounter{enumi}{4}
\item
  Ok, let's do some runs to make sure that we haven't changed anything
  that alters the ``no evolution'' default state of the model AND that
  when we allow gene flow (but nothing else), we get the expected
  result. First let's run two trials that should give us flat lines (we
  stay in HW over time), just to make sure everything is working
  properly:

  N

  p0

  toggledrift

  wA1A1

  wA1A2

  wA2A2

  pm

  m

  Description

  500

  0.1

  0

  1

  1

  1

  0

  0

  No evolution

  500

  0.1

  0

  1

  1

  1

  1

  0

  pm=1 but m=0, still no evolution

  Go ahead and use Simulation Control to test these two runs. All lines
  should be flat. You should select the variables p adj, pm, and
  selecteq to graph. The pm line will be at y=0 when pm=0. And the
  selecteq line should be at y = p0 (0.1), which is also where the padj
  line will be (uncheck the boxes below the graphs to reveal overlapping
  lines). You can also open the Control Panel and choose which datasets
  are available on the right. Note that in addition to graphs, you can
  make tables that will show you the value of your variable(s) for each
  time step. To do that, select the variables you want to include (use
  Shift to select multiple) and then click on one of the table icons,
  below the graph icon. I like to use the `Table Time Down' option. If
  you want to keep a table, you can save it as a text file if you click
  on File -\textgreater{} Save while you have the table selected. Or you
  can take a screenshot to keep the image. You can also copy and paste
  into a spreadsheet file and that will preserve the table grid
  (copy+paste into a text or word file will not keep the table grid).
\item
  Now, if that went ok we can start testing gene flow! Run the following
  parameter set but predict what will happen before you click Simulate
  and check the graph.

  N

  p0

  toggledrift

  wA1A1

  wA1A2

  wA2A2

  pm

  m

  Description

  500

  0.1

  0

  1

  1

  1

  1

  0.15

  15\% migration rate where all migrants are A1A1

  If you predicted the outcome correctly, great, proceed to work through
  the questions. If not, spend some time figuring out what's going on
  before you proceed - make sure you feel comfortable that you
  understand what to expect from gene flow.
\end{enumerate}

\subsection{Lab Write-Up}\label{lab-write-up}

Remember that the way we are representing gene flow, the values of pm
and m do not change over time, so the expected equilibrium for gene flow
(alone) will be that p adj = pm. Looking at the way in which we included
gene flow in the model, you might guess that m, the rate of migration,
will be important in determining how quickly we reach equilibrium. But
you might also notice that none of the changes we made involve N, so
equilibrium probably isn't going to depend on population size when only
gene flow is occurring. When we investigate how gene flow interacts with
selection and with drift, remember to identify which variables are the
most important for determining equilibrium for each of those forces
(fitness values and population size, respectively). For this lab
exercise, we have two primary goals: * Examine how gene flow alone
changes allele and genotype frequencies (with no selection and no
drift). Determine what factors have the most influence. * Investigate
the combined effects of gene flow and drift, gene flow and selection,
and all three forces at once. Explore under what conditions gene flow
alters the equilibrium outcome compared to when no gene flow was
present.

\subsubsection{Questions}\label{questions}

\begin{enumerate}
\def\labelenumi{\arabic{enumi}.}
\item
  Gene flow alone: Varying N, pm, and m. Try to predict what effect you
  think changing each of these variables will have on the outcome.
  Consider the equilibrium value of p and time to equilibrium. Run the
  conditions in the table below and compare the outcomes. Address the
  following questions in a paragraph-form summary. You may include
  graphs to support your conclusions.

  \begin{itemize}
  \tightlist
  \item
    Varying N: Compare Trials 1a and 1b. Did changing population size
    matter?\\
  \item
    Varying pm: Compare Trials 1a, 1c, and 1d. Did changing pm matter?\\
  \item
    Varying m: Compare Trials 1a, 1e, 1f, and 1g. Did changing m matter?

    Trial

    N

    p0

    toggledrift

    wA1A1

    wA1A2

    wA2A2

    pm

    m

    Description

    1a

    500

    0.1

    0

    1

    1

    1

    0.6

    0.15

    baseline, m = 0.15

    1b

    50

    0.1

    0

    1

    1

    1

    0.6

    0.15

    N = 50, smaller N

    1c

    500

    0.1

    0

    1

    1

    1

    0.9

    0.15

    pm = 0.9, higher pm

    1d

    500

    0.1

    0

    1

    1

    1

    0.3

    0.15

    pm = 0.3, lower pm

    1e

    500

    0.1

    0

    1

    1

    1

    0.6

    0.3

    m = 0.3, higher m

    1f

    500

    0.1

    0

    1

    1

    1

    0.6

    0.075

    m = 0.075, lower m

    1g

    500

    0.1

    0

    1

    1

    1

    0.6

    0.035

    m = 0.035, lower m
  \end{itemize}
\item
  Drift and gene flow. Keep in mind what you observed with only gene
  flow. Remember that for genetic drift, N is a crucial factor. Also,
  having very low allele frequencies leads to a higher probability of
  extinction so we'll consider what happens when an allele is rare in
  residents but common in immigrants. Remember drift will only influence
  our sub-population, not the source population of migrants; pm will be
  constant. Note that you will have to keep track of how many times you
  run each set of conditions. After observing the trials below, write a
  summary of the interaction between drift and gene flow. Does the
  presence of both forces prevent the equilibria you'd expect from only
  one of them acting? What does this illustrate for us about how genetic
  variation can be maintained in natural populations?

  Trial

  N

  p0

  toggledrift (\# runs)

  wA1A1

  wA1A2

  wA2A2

  pm

  m

  Description

  2a

  500

  0.1

  1-10

  1

  1

  1

  0.6

  0

  no migration, large N

  2b

  500

  0.1

  1-10

  1

  1

  1

  0.6

  0.035

  low migration, large N

  2c

  500

  0.1

  1-10

  1

  1

  1

  0.6

  0.1

  high migration, large N

  2d

  50

  0.1

  1-10

  1

  1

  1

  0.6

  0

  no migration, small N

  2e

  50

  0.1

  1-10

  1

  1

  1

  0.6

  0.035

  low migration, small N

  2f

  50

  0.1

  1-10

  1

  1

  1

  0.6

  0.1

  high migration, small N
\item
  Selection and gene flow. Now let's look at the outcome when both
  selection and gene flow are occuring in our sub-population (but no
  drift allowed). For selection, we remember that the difference in
  fitness between genotypes is important to determining the equilibrium
  and that the starting allele frequency and dominance coefficient
  influence how quickly we get to equilibrium. Right now, let's just
  focus on varying selection strength and the rate of migration to
  explore the intersection. We'll use a model of codominance of A1 and
  A2 (h = 0.5), with the A1A1 individual having the highest fitness.
  Complete the trials below and compare outcomes. Write a summary of the
  interaction between selection and gene flow.

  Trial

  N

  p0

  toggledrift (\# runs)

  wA1A1

  wA1A2

  wA2A2

  pm

  m

  Description

  3a

  500

  0.1

  0

  1

  0.9

  0.8

  0.6

  0

  no migration, s = 0.2

  3b

  500

  0.1

  0

  1

  0.9

  0.8

  0.6

  0.035

  low migration, s = 0.2

  3c

  500

  0.1

  0

  1

  0.9

  0.8

  0.6

  0.1

  high migration, s = 0.2

  3d

  500

  0.1

  0

  1

  0.75

  0.5

  0.6

  0.1

  high migration, s = 0.5
\item
  Selection, drift, and gene flow Now we put all three forces together!
  Since we just explored the intersection of selection and gene flow,
  we'll use some of those conditions and vary N to get a sense of how
  adding drift changes things. Run the following conditions, turning on
  drift and running 10 times each. Observe whether drift changes the
  outcome. Normally, we expect drift to lead to fixation or extinction
  of our A1 allele, at equilibrium. Does one allele ever achieve
  fixation or extinction in these runs? Summarize your observations on
  what happens to allele frequencies when all three evolutionary forces
  are acting simultaneously. Be sure to comment on whether equilibrium
  is reached and how it differs from forces acting alone. How do you
  think these simulations relate to the patterns we observe when we
  sample populations in nature?

  Trial

  N

  p0

  toggledrift (\# runs)

  wA1A1

  wA1A2

  wA2A2

  pm

  m

  Description

  4a

  500

  0.1

  1-10

  1

  0.9

  0.8

  0.6

  0.035

  N = 500

  4b

  100

  0.1

  1-10

  1

  0.9

  0.8

  0.6

  0.035

  N = 100

  4c

  50

  0.1

  1-10

  1

  0.9

  0.8

  0.6

  0.035

  N = 50

  4d

  25

  0.1

  1-10

  1

  0.9

  0.8

  0.6

  0.035

  N = 25
\item
  Reflect back on this lab exercise. In this exercise, we considered
  just one population and did not allow for variation over time in
  migration rate, population size, or selection strength. Imagine that
  we generalized these results to a situation with multiple populations
  and gene flow in any direction. How do you think the dynamics of
  allele frequency changes might be similar or different? What do you
  think are the takeaways that we should keep in mind when approaching
  conservation issues?
\end{enumerate}

\textbf{Remember to submit your model and responses via the
\href{https://forms.gle/E3QLACX1aA2Q6ist5}{Google Drive Form}.}

\begin{verbatim}
\end{verbatim}
  </main><aside class="col-md-3"><nav id="toc" aria-label="Table of contents"><h2>On this page</h2>
    </nav></aside></div>



    <footer><div class="pkgdown-footer-left">
  <p>Developed by Prof. A. Roles.</p>
</div>

<div class="pkgdown-footer-right">
  <p>Site built with <a href="https://pkgdown.r-lib.org/" class="external-link">pkgdown</a> 2.1.3.</p>
</div>

    </footer></div>





  </body></html>
